\documentclass[11pt]{article}

% ============================================================
% Packages and page layout
% ============================================================

\usepackage[T1]{fontenc}
\usepackage[utf8]{inputenc}
\usepackage[
    left=0.7in,
    right=0.7in,
    top=0.5in,
    bottom=0.5in
]{geometry}

\usepackage{xcolor}
\usepackage{array}
\usepackage{tabularx}
\usepackage{enumitem}
\usepackage{needspace}
\usepackage{hyperref}
\usepackage{qrcode}

\hypersetup{
    colorlinks=true,
    linkcolor=blue,
    urlcolor=blue,
    citecolor=blue,
    anchorcolor=blue
}

\urlstyle{same}
\pagenumbering{gobble}

\setlength{\parindent}{0pt}
\setlength{\parskip}{0pt}
\setlength{\emergencystretch}{2em}

\setlist[itemize]{
    leftmargin=1.4em,
    topsep=4pt,
    itemsep=3pt,
    parsep=2pt,
    partopsep=0pt
}

% ============================================================
% Reusable commands
% ============================================================

% Section heading
\newcommand{\cvsection}[1]{%
    \par\addvspace{14pt}%
    \Needspace{5\baselineskip}%
    {\large\bfseries #1\par}%
    \vspace{4pt}%
    \hrule
    \vspace{8pt}%
}

% Title on the left; date on the far right.
% Long titles wrap without overlapping the date.
\newcommand{\cvheading}[2]{%
    \par
    \Needspace{4\baselineskip}%
    \noindent
    \begin{tabularx}{\linewidth}{%
        @{}
        >{\raggedright\arraybackslash}X
        >{\raggedleft\arraybackslash}p{1.65in}
        @{}
    }
        \textbf{#1} & #2
    \end{tabularx}\par
}

% Entry with title, date, description, and itemize contents.
% Usage:
% \cventry{Title}{Date}{Description}{\item ... \item ...}
\newcommand{\cventry}[4]{%
    \cvheading{#1}{#2}%
    \if\relax\detokenize{#3}\relax
    \else
        \smallskip
        #3\par
    \fi
    \begin{itemize}
        #4
    \end{itemize}
    \addvspace{6pt}%
}

% Entry without bullet points.
% Usage: \cvplainentry{Title}{Date}{Description}
\newcommand{\cvplainentry}[3]{%
    \cvheading{#1}{#2}%
    \if\relax\detokenize{#3}\relax
    \else
        \smallskip
        #3\par
    \fi
    \addvspace{8pt}%
}

% Study-topic category and contents
\newcommand{\StudyTopics}[2]{%
    \par
    \begin{tabularx}{\linewidth}{%
        @{}
        >{\raggedright\arraybackslash}p{0.18\linewidth}
        >{\raggedright\arraybackslash}X
        @{}
    }
        \textbf{#1} & #2
    \end{tabularx}\par
    \addvspace{6pt}%
}

% ============================================================
% Document
% ============================================================

\begin{document}

% ============================================================
% Contact information
% ============================================================

\noindent
\begin{minipage}[c]{0.15\textwidth}
    \hfill
\end{minipage}%
\begin{minipage}[c]{0.70\textwidth}
    \centering
    {\LARGE\bfseries Dovan Darazi\par}
    \vspace{5pt}
    Beirut-Lebanon $|$ (+961) 78 896712\par
    \vspace{3pt}
    \href{mailto:dovan_darazi@outlook.com}
         {\textcolor{black}{dovan\_darazi@outlook.com}}\par
\end{minipage}%
\begin{minipage}[c]{0.15\textwidth}
    \raggedleft
    \qrcode[height=1.5cm]{https://dovandarazi.github.io/}
\end{minipage}

% ============================================================
% Research interests
% ============================================================

\cvsection{RESEARCH INTERESTS}

High Energy Physics: Theory and Computational/Experimental,
Astroparticle Physics, Early Universe Physics, Dark Matter
Physics, Cosmology, General Relativity.

% ============================================================
% Academic history
% ============================================================

\cvsection{ACADEMIC HISTORY}

\cvplainentry
{American University of Beirut}
{Jan 2025 -- Present}
{%
    \textbf{Masters of Science in Physics}\par
    \smallskip
    Electives: Quantum Field Theory, Advanced Mathematical Methods
}

\cvplainentry
{Lebanese University, Faculty of Sciences II}
{Oct 2021 -- Jul 2024}
{%
    \textbf{Bachelor of Science in Physics}\par
    \smallskip
    \textbf{GPA:} 3.6 (passing grade 60\%)\par
    \textbf{Rank:} first year: 1/17; second year: 1/3;
    third year: 1/1\par
    \smallskip
    Elective courses: Waves and vibrations, Fluid mechanics,
    Classical Field Theory, Laser Physics, Thermal Transfer,
    Material Science, Molecular Physics, Medical Imaging
}

% ============================================================
% Research experience
% ============================================================

\cvsection{RESEARCH EXPERIENCE}
\cventry
{Master's Thesis: Investigating Baryogenesis through B-Meson Decays}
{May 2026 -- Present}
{American University of Beirut --- BaBar data analysis with Dr. Racha Cheaib}
{
    \item Investigating a possible mechanism for baryogenesis through B-meson decays involving a $\Xi_c^{+}$ baryon and a dark matter particle.
    \item Developing an analysis using BaBar data to search for this decay signature and explore its potential connection to the matter--antimatter asymmetry of the Universe.
}


\cventry
{Research assistant, American University of Beirut (AUB)}
{Jun 2025 -- Oct 2025}
{Laboratory for Plasma and Fluid Dynamics}
{%
    \item Designed, built, and tested an improved divertor model
    for magnetic confinement nuclear fusion devices aimed at
    filtering high-energy debris within the plasma. The design
    replaces conventional solid components with a flowing liquid
    metal system capable of self-healing after repeated particle
    impacts, significantly reducing material degradation and
    maintenance requirements. The system captures and removes
    debris from the plasma, filters it through the liquid medium,
    and recirculates the purified material back into the divertor,
    enhancing durability and operational efficiency.
}

\cventry
{Research assistant, American University of Beirut (AUB)}
{Sep 2023 -- Oct 2025}
{Laboratory for Plasma and Fluid Dynamics}
{%
    \item Project lead on developing a high-speed stereoscopic
    diagnostics device in order to detect and identify UFOs,
    Flakes and Dust particles in the plasma flow of a magnetic
    confinement nuclear fusion reactor. The objective of this
    project is to determine particles' sizes as well as their
    positions in three dimensions and their velocities in
    component form. This data plays a pivotal role in informing
    the necessary adjustments to the input parameters of the
    device, enabling us to anticipate and mitigate potential
    disruptions to the plasma flow with greater precision.

    \item \textbf{Access The Project Documents:}
    \href{https://dovandarazi.github.io/\#professional-experience}
         {\textbf{here}}
}

\cventry
{Intern, Lebanese American University (LAU)}
{Jul 2023 -- Sep 2023}
{}
{%
    \item Assisted in data collection, organization, and analysis
    for a Research project on the water-energy-food nexus studying
    the impact of the Lebanese economic and energy crisis on
    water security. Worked on finalizing the paper for publication.
}

% ============================================================
% Publications and presentations
% ============================================================

\cvsection{PUBLICATIONS AND PRESENTATIONS}

\textbf{Publications:}
\href{https://scholar.google.com/citations?hl=en&user=SUrLchwAAAAJ}
     {Google Scholar}

\begin{itemize}
    \item
    \href{https://doi.org/10.2166/aqua.2024.329}
         {Navigating the water--energy nexus amidst the Lebanese
          economic crisis}
\end{itemize}

\smallskip
\textbf{Presentations}

\begin{itemize}
    \item Presented my project on imaging disruptions in magnetic
    confinement fusion devices at the first MENA region plasma
    fusion physics (MPFu) conference in Gammarth, Tunisia, as the
    youngest presenter.
\end{itemize}

% ============================================================
% Academic awards
% ============================================================

\cvsection{ACADEMIC AWARDS}

\cvplainentry
{Graduate Fellowship and Assistantship Program}
{Jan 2025 -- Present}
{American University of Beirut --- Beirut, Lebanon}

% ============================================================
% Teaching and mentoring experience
% ============================================================

\cvsection{TEACHING AND MENTORING EXPERIENCE}

\cventry
{Teaching Assistant, American University of Beirut (AUB)}
{Jan 2025 -- Present}
{Teaching Physics labs for engineers and premeds
(PHYS 211L \& PHYS 206L)}
{%
    \item My duties as a TA were: Deliver a lecture on how to do
    the experiment, guide and assist the students while they're
    doing the experiment, correct reports that the students submit
    for each experiment, Design and correct quizzes that test the
    students' preparation, have office hours for followups,
    questions, and inquiry, correct final exams along with other
    teaching assistants.
}

\cventry
{Particle Physics Summer School Instructor and Mentor}
{Jul 2026 -- Aug 2026}
{Summer Research Experience Program (SREP): Analyzing LHCb Data}
{%
    \item Mentored a group of high school students, teaching the
    fundamentals of particle physics, detector physics, Python
    programming, and large-scale data analysis.

    \item Guided students in analyzing LHCb B-meson decay data,
    including particle identification and developing code to
    reconstruct specified decay channels using charge and
    four-momentum conservation.

    \item Supervised the selection and optimization of analysis
    cuts by evaluating the trade-off between efficiency and purity,
    culminating in a measurement of the asymmetry between B$^{+}$
    and B$^{-}$ decays.

    \item Supported students in presenting their methods and
    findings to fellow participants and physics faculty at the
    conclusion of the program.
}

\cventry
{Particle Physics Summer School Instructor and Mentor}
{Jul 2026 -- Aug 2026}
{Summer Research Experience Program (SREP):
Designing a Dark Photon Detector}
{%
    \item Mentored a group of high school students, teaching the
    fundamentals of particle physics, detector physics, Python
    programming, and large-scale data analysis.

    \item Guided students in exploring different types of particle
    detectors and their operating principles, applying this
    knowledge to develop a conceptual detector design for a
    research objective of their choice.

    \item Supervised the team's development of a conceptual
    detector designed to search for dark photons.
}

% ============================================================
% Technical skills
% ============================================================

\cvsection{TECHNICAL SKILLS}

\noindent
\begin{tabularx}{\linewidth}{%
    @{}
    >{\bfseries\raggedright\arraybackslash}X
    l
    >{\bfseries\raggedright\arraybackslash}X
    l
    @{}
}
    Python & Advanced
        & \LaTeX & Intermediate \\[4pt]
    ROOT & Intermediate
        & OriginLab & Intermediate \\[4pt]
    C++ & Intermediate
        & AutoCAD & Intermediate \\[4pt]
    Rust & Intermediate
        & Inventor & Intermediate \\[4pt]
    MATLAB & Intermediate
        & Blender & Beginner \\[4pt]
    Linux & Intermediate
        & COMSOL Multiphysics & Beginner \\[4pt]
    SQL & Beginner
        & Raspberry Pi & Beginner
\end{tabularx}

% ============================================================
% Projects
% ============================================================

\cvsection{PROJECTS}

\cventry
{Monte Carlo Simulation and Analysis of the Ising Model}
{Apr 2025 -- Jun 2025}
{Statistical Mechanics --- Final Course Project}
{
    \item Studied the Ising model and implemented a Monte Carlo simulation to investigate its statistical mechanical behavior.
    \item Analyzed the simulation results alongside the theoretical description of the model as part of the course's final project.
}


\cventry
{Investigating High-Efficiency Thermoelectric Metals
for Integration with Hollow Blocks}
{Dec 2024 -- Jan 2025}
{Numerical Analysis II course, Lebanese University,
Faculty of Sciences II, Department of Physics}
{%
    \item In this project, we aim to design an innovative system
    for sustainable energy generation by leveraging the temperature
    gradient between the interior and exterior of buildings.
    Working as a team of two, we are researching thermoelectric
    metals with high Seebeck coefficients for integration into
    hollow blocks (bricks). Our goal is to identify materials that
    are not only efficient in converting heat to electricity but
    are also cost-effective and abundant, ensuring scalability
    and accessibility for widespread use.

    \par\smallskip
    To achieve this, we are utilizing COMSOL Multiphysics to create
    detailed simulations of the proposed device, enabling us to
    model and optimize its performance under real-world conditions.
    By combining advanced material analysis with cutting-edge
    simulation tools, this project has the potential to pave the
    way for practical applications of thermoelectric technology in
    building materials, contributing to energy-efficient
    infrastructure and renewable energy solutions.
}


\cventry
{GOLEM Tokamak: Toroidal Field Experiment}
{Summer 2024}
{Virtually operated GOLEM tokamak, Czech Technical University in Prague}
{
    \item Investigated the impact of toroidal magnetic field strength on plasma stability and confinement by varying the toroidal field capacitor voltage.
    \item Studied its effects on key plasma parameters, including plasma current, edge safety factor, discharge duration, and electron energy confinement time.
}

\cventry
{Developing an Interactive Virtual Assistant Display}
{Mar 2024 -- Jun 2024}
{Open Laboratory, Lebanese University, Faculty of Sciences II,
Department of Physics}
{%
    \item Developing a 3D interface for the laboratory's virtual
    assistant using Blender, with the goal of enhancing user
    interaction and visualization. The interface is designed to
    be displayed on an anamorphic screen, providing a dynamic
    and immersive visual experience. This project serves as a
    foundation for creating a more advanced projection system,
    aimed at improving display technologies for laboratory
    applications and beyond.
}

\cventry
{Studying the Barber pole effect in sucrose solution}
{Oct 2023 -- Feb 2024}
{Optics Laboratory Project, Lebanese University,
Faculty of Sciences II}
{%
    \item As part of a student project supervised by a professor
    at the Laboratory of Biomedical Imaging (LBMI), Lebanese
    University, Faculty of Sciences II, our team conducted a series
    of experiments to investigate the barber pole effect. The
    project involved designing and performing experiments to
    observe the phenomenon and collect relevant experimental data.

    \par\smallskip
    Through comprehensive data analysis, we validated the
    theoretical concepts identified during the literature review,
    demonstrating alignment between our findings and existing
    research. The results were presented to the supervising
    professor and peers, showcasing our approach, methodology,
    and key insights. This project provided valuable hands-on
    experience in experimental design, data analysis, and
    effective scientific communication.
}

\cventry
{Studying and Using a Prism Spectrometer Goniometer}
{Oct 2023 -- Jan 2024}
{Optics Laboratory Course, Lebanese University,
Faculty of Sciences II}
{%
    \item Calibrated a non-digital spectrometer goniometer using a
    mercury lamp to ensure accuracy in subsequent measurements.
    The calibrated setup was utilized to determine the Balmer
    spectra of a hydrogen lamp and calculate the Rydberg constant.

    \par\smallskip
    A theoretical geometrical optical analysis was conducted to
    achieve the best fit for the measured data, with results
    visualized using Python's matplotlib library. To further
    enhance understanding, the studied concept was modeled in
    Blender, providing a 3D representation of the optical system.
    This project integrated experimental techniques, computational
    analysis, and 3D modeling to deepen insights into
    spectroscopic phenomena.
}

\cventry
{The Effects of Photon Energy and Magnetic Fields
on the Germination/Growth of Plants}
{Aug 2023 -- Oct 2023}
{Student-led project under the supervision of a professor
(LBMI, Lebanese University, Faculty of Sciences II)}
{%
    \item Designed and optimized experimental methods for seed
    germination and growth under controlled conditions. The
    experiments investigated the effects of blue and red light
    (individually) on germination, as well as the influence of a
    magnetic field on the process.

    \par\smallskip
    The study aimed to explore the potential role of quantum
    processes in seed germination and growth, providing insights
    into how environmental factors such as light wavelength and
    magnetic fields impact these biological mechanisms. This
    project combined experimental precision with theoretical
    exploration to advance understanding in biophysics and
    plant sciences.
}

% ============================================================
% Workshops, conferences, and schools
% ============================================================

\cvsection{WORKSHOPS, CONFERENCES, AND SCHOOLS}

\cvplainentry
{Physics Without Frontiers:
Lebanon's First Quantum Technology Workshop}
{Oct 2025}
{Abdus Salam International Centre for Theoretical Physics}

\cvplainentry
{STARS: An Introduction to their Structure and Evolution}
{Sep 2025}
{}

\cvplainentry
{The fourth MPFuS}
{Summer 2024}
{Attended the fourth MPFuS summer school for plasma and fusion
physics held in Gammarth, Tunisia.}

% ============================================================
% Studied topics
% ============================================================

\cvsection{STUDIED TOPICS}

\StudyTopics
{Physics}
{Particle Physics, Quantum Mechanics, Quantum Field Theory,
Nuclear Physics, Classical Field Theory, General Relativity,
Computational Physics, Plasma Physics, Fluid Mechanics,
Experimental Physics, Electric Circuits, Error Analysis,
Classical Mechanics, Statistical Mechanics, Electromagnetism,
Waves And Vibrations, Laser Physics, Thermal Transfer,
Material Science, Molecular Physics, Medical Imaging Physics.}

\StudyTopics
{Mathematics}
{Tensor Analysis, Linear Algebra, Differential Geometry,
Topology, Group Theory, Complex Analysis, Fourier Series and
Fourier Transform, Laplace Transform, Ordinary and Partial
Differential Equations, Numerical Computing and Analysis,
Real Analysis}

\StudyTopics
{Space Science}
{Astronomy, Astrophysics, Planetary Geology}

% ============================================================
% Online courses
% ============================================================

\cvsection{ONLINE COURSES}

\textbf{\href{https://www.coursera.org/}{coursera.org}}

\smallskip
Particle Physics: An Introduction, From the Big Bang to Dark
Energy, Astro 101: Black Holes, Astronomy: Exploring Time and
Space, Astrobiology: Exploring Other Worlds, The Science of
the Solar System, How to Write and Publish a Scientific Paper

% Retained from the original source as commented-out content:
% Supplemental studies (books):
% Introduction to Electrodynamics by Griffith,
% Tensor calculus for physics,
% A first course in general relativity 2E

% ============================================================
% Certifications
% ============================================================

\cvsection{CERTIFICATIONS}

\begin{itemize}
    \item Raspberry Pi Certification from IEEE
    \item First Aid Certification from Red Cross
\end{itemize}

% ============================================================
% Languages
% ============================================================

\cvsection{LANGUAGES}

\begin{tabularx}{\linewidth}{%
    @{}
    >{\bfseries}l
    X
    @{}
}
    English & Fluent \\[4pt]
    Arabic & Native \\[4pt]
    French & Conversational
\end{tabularx}


\end{document}